\documentclass[11pt, a4paper]{article}

%% Packages
\usepackage[margin=2.5cm]{geometry}
\usepackage{amsmath, amssymb}
\usepackage{booktabs}
\usepackage{array}
\usepackage{hyperref}
\usepackage{xcolor}
\usepackage{graphicx}
\usepackage{caption}
\usepackage{subcaption}
\usepackage{parskip}
\usepackage{microtype}
\usepackage{enumitem}
\usepackage{fancyhdr}
\usepackage{titlesec}
\usepackage{listings}
\usepackage{nicematrix}

%% Hyperref config
\hypersetup{
  colorlinks=true,
  linkcolor=blue!60!black,
  urlcolor=blue!60!black,
  citecolor=blue!60!black,
}

%% Header/footer
\pagestyle{fancy}
\fancyhf{}
\fancyhead[L]{\small BTCUSDT Perpetual Futures: Short-Horizon Signal Research}
\fancyhead[R]{\small \thepage}
\renewcommand{\headrulewidth}{0.4pt}

%% Section formatting
\titleformat{\section}{\large\bfseries}{{\thesection.}}{0.6em}{}[\titlerule]
\titleformat{\subsection}{\normalsize\bfseries}{{\thesubsection}}{0.5em}{}

%% Table column type
\newcolumntype{C}[1]{>{\centering\arraybackslash}p{#1}}

%% Listings style (for inline code)
\lstset{
  basicstyle=\ttfamily\small,
  backgroundcolor=\color{gray!10},
  frame=single,
  framesep=3pt,
  rulecolor=\color{gray!40},
  breaklines=true,
}

%% -----------------------------------------------------------------------
\begin{document}

\begin{titlepage}
  \centering
  \vspace*{2cm}
  {\LARGE\bfseries Short-Horizon Mean-Reversion Signals\\[0.4em]
   for BTCUSDT Perpetual Futures\par}
  \vspace{1cm}
  {\large Quantitative Research Report\par}
  \vspace{0.4cm}
  {\normalsize Tick-to-OHLCV Pipeline $\cdot$ 8-Signal Composite
               $\cdot$ Cost-Aware Vectorized Backtest\par}
  \vspace{1.5cm}
  \tableofcontents
  \vfill
  {\normalsize March 2026\par}
\end{titlepage}

%% -----------------------------------------------------------------------
\section{Introduction}

This report documents an end-to-end quantitative research workflow for developing
short-horizon directional signals on Binance BTCUSDT perpetual futures.
The study targets a \textbf{5--15 minute holding period}, a regime that sits at the
boundary between market-microstructure phenomena and the slower price discovery
captured by conventional technical strategies.

Starting from two days of raw, nanosecond-resolution trade-level data (approximately
2.3~million events), the work proceeds through four stages:

\begin{enumerate}[leftmargin=2em]
  \item \textbf{Data exploration and OHLCV aggregation} --- data-quality audits,
        distributional characterisation of 1-minute returns, autocorrelation
        analysis, order-flow decomposition, and intraday seasonality profiling.
  \item \textbf{Signal design} --- eight individual signals spanning mean-reversion,
        order-flow, and momentum families, combined into a weighted composite
        with exponential smoothing and hysteresis-based position discretisation.
  \item \textbf{Backtesting} --- a vectorized simulation framework with realistic
        execution assumptions, transaction-cost sensitivity sweeps, rebalance-frequency
        analysis, and a walk-forward out-of-sample validation.
  \item \textbf{Critical analysis} --- an honest evaluation of what the results
        do and do not prove, including identifiable biases and the conditions
        required for live profitability.
\end{enumerate}

The codebase is organised as a Python package (\texttt{src/}) with three
companion Jupyter notebooks:

\begin{itemize}[leftmargin=2em, itemsep=0pt]
  \item \texttt{notebooks/01\_data\_exploration.ipynb}
  \item \texttt{notebooks/02\_signal\_research.ipynb}
  \item \texttt{notebooks/03\_backtest.ipynb}
\end{itemize}

%% -----------------------------------------------------------------------
\section{Data}

\subsection{Source and Schema}

The dataset comprises two Parquet files covering UTC days
\textbf{18 April 2025} and \textbf{19 April 2025}.
Each row records a single aggressed trade on Binance BTCUSDT perpetual futures:

\begin{center}
\begin{tabular}{lll}
\toprule
\textbf{Field} & \textbf{Type} & \textbf{Description} \\
\midrule
\texttt{trade\_time}  & int64 (ns)   & Exchange-assigned nanosecond timestamp \\
\texttt{side}         & \{BID, ASK\} & BID = buyer-initiated; ASK = seller-initiated \\
\texttt{price}        & float64      & Execution price (USDT) \\
\texttt{quantity}     & float64      & Traded quantity (BTC) \\
\bottomrule
\end{tabular}
\end{center}

Two derived columns are added during loading:
\texttt{signed\_qty} (positive for BID, negative for ASK) and
\texttt{notional} ($= \mathrm{price} \times \mathrm{quantity}$).

\subsection{Coverage and Quality}

\begin{center}
\begin{tabular}{lr}
\toprule
\textbf{Statistic} & \textbf{Value} \\
\midrule
Total trade records      & 2{,}308{,}624 \\
Price range              & \$84{,}260 -- \$85{,}600 \\
Missing values (any field) & 0 \\
Minutes with zero trades & 0 / 2{,}880 \\
Mean trades per minute   & 802 \\
Max trades per minute    & 14{,}663 \\
BID trades               & 1{,}144{,}581 (49.6\%) \\
ASK trades               & 1{,}164{,}043 (50.4\%) \\
Median trade size (BTC)  & 0.005 \\
\bottomrule
\end{tabular}
\end{center}

The data are \textit{dense} (no gaps), clean (zero nulls), and well-balanced
between buyer- and seller-initiated flow.  Duplicate nanosecond timestamps
(1{,}504{,}842 rows share a timestamp with at least one other row) reflect
batched matching-engine outputs at the same microsecond, which is normal
for a high-throughput exchange.

\subsection{OHLCV Aggregation}

Tick data are resampled into \textbf{1-minute bars} using a standard
open/high/low/close/volume scheme, extended with microstructure fields:

\[
\mathrm{VWAP}_t
  = \frac{\sum_{i \in t} p_i q_i}{\sum_{i \in t} q_i},
\qquad
\mathrm{buy\_ratio}_t
  = \frac{V^{\mathrm{buy}}_t}{V^{\mathrm{buy}}_t + V^{\mathrm{sell}}_t},
\]

\noindent where $V^{\mathrm{buy}}_t$ and $V^{\mathrm{sell}}_t$ are the
buyer- and seller-initiated volumes in bar $t$.  In addition:

\begin{itemize}[leftmargin=2em, itemsep=0pt]
  \item \textbf{log returns}: $r_t = \ln(C_t / C_{t-1})$
  \item \textbf{spread}: $h_t - l_t$ (high minus low within the bar)
  \item \textbf{trade count}: number of individual fills per bar
\end{itemize}

A total of \textbf{2{,}880 one-minute bars} are produced across the two-day window.

\subsection{Return Distribution}

One-minute log returns exhibit the distributional properties typical of
liquid crypto markets:

\begin{center}
\begin{tabular}{lr}
\toprule
\textbf{Statistic} & \textbf{Value} \\
\midrule
Mean   & $\approx 0$ \\
Std    & $\approx 0.00015$ (1.5 bps) \\
Skewness  & small, close to zero \\
Excess kurtosis & $> 3$ (heavy tails) \\
\bottomrule
\end{tabular}
\end{center}

The QQ plot against the normal distribution confirms pronounced heavy tails ---
the distribution is leptokurtic.  This motivates signal normalization via rolling
z-scores (rather than raw levels) and robust clipping to $[-3, 3]$ to limit
the influence of outlier bars.

\subsection{Autocorrelation: Motivation for Mean Reversion}

A critical finding is the \textbf{negative lag-1 autocorrelation} of 1-minute
log returns:

\begin{center}
\begin{tabular}{cc}
\toprule
\textbf{Lag (bars)} & \textbf{Autocorrelation} \\
\midrule
1 & $-0.05$ to $-0.08$ \\
2 & small, mixed \\
3--5 & near zero \\
\bottomrule
\end{tabular}
\end{center}

A down move at the 1-minute bar level is slightly more likely to be followed
by an up move than another down move.  This pattern is consistent with
\textit{bid-ask bounce} and \textit{microstructure mean reversion}: market
makers systematically quote on both sides and tighten spreads after one-sided
flow.  It is the primary empirical motivation for the mean-reversion signal
suite developed in Section~\ref{sec:signals}.

\subsection{Intraday Patterns}

Hourly breakdowns of volume, volatility, and trade count reveal:

\begin{itemize}[leftmargin=2em, itemsep=0pt]
  \item Elevated activity in the Asian morning session (roughly 00:00--04:00 UTC)
        and the US afternoon session (roughly 13:00--17:00 UTC).
  \item Return volatility is positively correlated with trade count
        (more activity = wider intrabar ranges).
  \item Buy ratio oscillates around 0.50 with no persistent directional bias,
        confirming the absence of a multi-hour trending regime in this sample.
\end{itemize}

%% -----------------------------------------------------------------------
\section{Signal Design}
\label{sec:signals}

\subsection{Design Philosophy}

All eight signals are designed around a single empirical observation: at horizons
of 1--15 minutes, BTCUSDT returns exhibit \textit{mean reversion} rather than
momentum.  Seven of the eight signals generate a \textbf{positive (long) reading
when price or flow is depressed relative to a recent reference level}, and vice versa.
The eighth signal (trade intensity) captures abnormal activity and pairs it with
the contemporaneous price direction.

Signal convention throughout:
\[
s_t > 0 \implies \text{go long}; \quad s_t < 0 \implies \text{go short.}
\]

\subsection{Individual Signals}

\paragraph{Signal 1 --- Z-Score Mean Reversion.}
\[
z_t = \frac{C_t - \bar{C}_{t,L}}{\sigma_{C,t,L}},
\qquad
s^{(1)}_t = -z_t,
\]
where $\bar{C}_{t,L}$ and $\sigma_{C,t,L}$ are the rolling mean and standard
deviation of close prices over a lookback of $L = 20$ bars.  A below-average
price (negative $z$) generates a positive signal.

\paragraph{Signal 2 --- Bollinger Band Mean Reversion.}
\[
s^{(2)}_t = -\operatorname{clip}\!\left(
  \frac{C_t - \bar{C}_{t,L}}{k \cdot \sigma_{C,t,L}},\ -1,\ 1
\right),
\quad L=20,\ k=2.
\]
This normalizes the position within the Bollinger envelope to $[-1, 1]$ and
inverts for mean reversion.  Signals 1 and 2 are algebraically very similar
when the clip is not binding; their correlation in practice approaches unity
(see Section~\ref{sec:discussion}).

\paragraph{Signal 3 --- VWAP Reversion.}
\[
\mathrm{VWAP}^{(L)}_t
  = \frac{\sum_{\tau=t-L+1}^{t} N_\tau}{\sum_{\tau=t-L+1}^{t} V_\tau},
\qquad
s^{(3)}_t
  = -\operatorname{clip}\!\left(
     \frac{(C_t - \mathrm{VWAP}^{(L)}_t) / \mathrm{VWAP}^{(L)}_t}
          {\sigma_{\mathrm{dev},t,L}},\ -3,\ 3
   \right), \quad L=15.
\]
Here $N_\tau$ is notional and $V_\tau$ is volume.  VWAP acts as a
\textit{volume-weighted fair value} anchor, which is less susceptible to
isolated large price moves than a simple moving average.

\paragraph{Signal 4 --- Order Flow Imbalance (OFI).}
\[
\phi_t = \frac{V^{\mathrm{buy}}_t - V^{\mathrm{sell}}_t}{V_t},
\qquad
s^{(4)}_t
  = -\operatorname{clip}\!\left(
     \frac{\bar{\phi}^{\mathrm{fast}}_t - \bar{\phi}^{\mathrm{slow}}_t}
          {\sigma_{\Delta\phi,t}},\ -3,\ 3
   \right), \quad f=5,\ s=20.
\]
The fast-minus-slow decomposition isolates \textit{recent} imbalance above
the baseline.  Excess buying ($\phi > 0$) is faded contrarily: informed
order flow pushes price in the short run, but the resulting
overextension mean-reverts.

\paragraph{Signal 5 --- Trade Intensity.}
\[
I_t = \frac{n_t - \bar{n}_{t,L}}{\sigma_{n,t,L}},
\qquad
s^{(5)}_t
  = \operatorname{clip}\!\left(I_t \cdot \operatorname{sgn}(r_t),\ -3,\ 3
  \right), \quad L=20.
\]
This combines an activity z-score with the sign of the contemporaneous return:
a volume spike accompanied by an up move is treated as a \textit{long} signal
(informed aggressive buyers).  Unlike the other signals, this is a
\textbf{momentum} signal, not contrarian.

\paragraph{Signal 6 --- RSI Mean Reversion.}
\[
\mathrm{RSI}_t = 100 - \frac{100}{1 + \mathrm{RS}_t},
\quad
\mathrm{RS}_t = \frac{\overline{\Delta^+}_{t,P}}{\overline{\Delta^-}_{t,P}},
\quad
s^{(6)}_t
  = \operatorname{clip}\!\left(-\frac{\mathrm{RSI}_t - 50}{20},\ -1.5,\ 1.5
  \right),
\quad P=14.
\]
Gains and losses are smoothed with an exponential weighted moving average (span
$= P$).  The signal is centred at 50 (neutral) and maps overbought (RSI~$>70$)
to $-1$ and oversold (RSI~$<30$) to $+1$.

\paragraph{Signal 7 --- MA Crossover (Contrarian).}
\[
s^{(7)}_t
  = -\operatorname{clip}\!\left(
     \frac{\mathrm{EMA}^{f}_t - \mathrm{EMA}^{s}_t}{\sigma_{C,t,s}},
     \ -3,\ 3
   \right),
\quad f=5,\ s=15.
\]
Short-term momentum (fast EMA above slow EMA) is faded, consistent with
the mean-reverting microstructure.

\paragraph{Signal 8 --- Volume-Weighted Momentum (Contrarian).}
\[
m_t = r_t \cdot \frac{V_t}{\bar{V}_{t,L}},
\qquad
s^{(8)}_t
  = -\operatorname{clip}\!\left(
     \frac{\sum_{\tau=t-L+1}^{t} m_\tau}{\sigma_{\Sigma m, t, 2L}},
     \ -3,\ 3
   \right),
\quad L=10.
\]
High-conviction up moves (large volume + positive return, cumulated over
$L$ bars) are faded.

\subsection{Composite Signal and Position Sizing}
\label{sec:composite}

All eight signals are combined into a weighted composite:
\[
C^{\mathrm{raw}}_t = \sum_{k=1}^{8} w_k s^{(k)}_t,
\]
with hand-specified weights that overweight the VWAP reversion signal on
account of its grounding in volume-fair-value:

\begin{center}
\begin{tabular}{lc}
\toprule
\textbf{Signal} & \textbf{Weight} $w_k$ \\
\midrule
Z-score MR        & 0.15 \\
Bollinger MR      & 0.15 \\
VWAP Reversion    & 0.20 \\
Order Flow Imb.   & 0.15 \\
Trade Intensity   & 0.05 \\
RSI MR            & 0.10 \\
MA Crossover      & 0.10 \\
Volume Momentum   & 0.10 \\
\midrule
Total             & 1.00 \\
\bottomrule
\end{tabular}
\end{center}

The composite is then smoothed with a 5-bar EMA to reduce whipsaw:
\[
C_t = \mathrm{EMA}_5\!\left(C^{\mathrm{raw}}_t\right).
\]

\paragraph{Hysteresis Position Sizing.}
Continuous signals are discretised into $\{-1, 0, +1\}$ via a hysteresis rule
with entry threshold $\theta_{\mathrm{in}} = 0.15$ and exit threshold
$\theta_{\mathrm{out}} = 0.05$:

\[
\mathrm{pos}_t =
\begin{cases}
+1 & \text{if } \mathrm{pos}_{t-1} = 0 \text{ and } C_t > \theta_{\mathrm{in}} \\
-1 & \text{if } \mathrm{pos}_{t-1} = 0 \text{ and } C_t < -\theta_{\mathrm{in}} \\
0  & \text{if } \mathrm{pos}_{t-1} = +1 \text{ and } C_t < \theta_{\mathrm{out}} \\
0  & \text{if } \mathrm{pos}_{t-1} = -1 \text{ and } C_t > -\theta_{\mathrm{out}} \\
\mathrm{pos}_{t-1} & \text{otherwise}
\end{cases}
\]

This \textit{gap} between entry and exit thresholds prevents the strategy from
flipping back and forth on noise.  In the data, this produces an average
holding period of approximately \textbf{14.7 bars ($\approx$15 minutes)},
which sits at the upper edge of the target regime.

%% -----------------------------------------------------------------------
\section{Backtesting Framework}

\subsection{Execution Assumptions}

\begin{itemize}[leftmargin=2em, itemsep=0pt]
  \item \textbf{Entry price}: close of the bar \textit{following} signal generation
        (one-bar delay via \texttt{signal.shift(1)}), eliminating look-ahead bias.
  \item \textbf{Return capture}: the strategy earns the close-to-close return of
        the entry bar.
  \item \textbf{Transaction cost model}:
        $\mathrm{cost}_t = |\Delta \mathrm{pos}_t| \times c$,
        where $c$ is the one-way cost in basis points.
  \item \textbf{No leverage constraints} or margin calls are modeled.
  \item \textbf{No market impact} or queue priority effects.
\end{itemize}

\subsection{PnL Decomposition}

Bar-level gross PnL:
\[
\pi^{\mathrm{gross}}_t = \mathrm{pos}_t \times r^{\mathrm{fwd}}_t,
\]
where $r^{\mathrm{fwd}}_t$ is the forward close-to-close return.

Net PnL:
\[
\pi^{\mathrm{net}}_t
  = \pi^{\mathrm{gross}}_t - |\Delta \mathrm{pos}_t| \times \frac{c}{10{,}000}.
\]

\subsection{Walk-Forward Validation}

To test out-of-sample generalization, a rolling walk-forward scheme is applied:
a \textbf{12-hour training window} (720 bars) is used to warm up the rolling
signal statistics, followed by a \textbf{4-hour test window} (240 bars) on
which PnL is recorded.  The window rolls forward by one test period and repeats.

Note: signal \textit{parameters} (lookback lengths, weights) are not
re-estimated in each training fold --- they remain fixed.  The training window
serves purely to initialise the rolling statistics (means, standard deviations)
that the signals depend on.  This is discussed critically in
Section~\ref{sec:discussion}.

%% -----------------------------------------------------------------------
\section{Results}

\subsection{Information Coefficient Analysis}

The Spearman rank correlation (IC) between each signal and $h$-bar
forward returns is computed across all bars:

\begin{center}
\begin{tabular}{lrrrr}
\toprule
\textbf{Signal} & \textbf{1-min IC} & \textbf{5-min IC}
  & \textbf{10-min IC} & \textbf{15-min IC} \\
\midrule
Z-score MR       & 0.046 & 0.097 & 0.092 & 0.085 \\
Bollinger MR     & 0.046 & 0.097 & 0.092 & 0.085 \\
VWAP Reversion   & 0.040 & 0.097 & 0.085 & 0.077 \\
OFI              & 0.052 & 0.107 & 0.092 & 0.074 \\
Trade Intensity  & 0.011 & 0.011 & 0.002 & \textbf{$-0.018$} \\
RSI MR           & 0.056 & 0.116 & 0.110 & 0.102 \\
MA Crossover     & 0.065 & 0.104 & 0.106 & 0.104 \\
Volume Momentum  & 0.043 & 0.078 & 0.108 & 0.077 \\
\midrule
\textbf{Composite}& \textbf{0.063} & \textbf{0.105} & \textbf{0.099} & \textbf{0.079} \\
\bottomrule
\end{tabular}
\end{center}

Key observations:
\begin{itemize}[leftmargin=2em, itemsep=0pt]
  \item All signals except Trade Intensity show positive IC across the 1--15 min range.
  \item IC peaks around the 5--10 minute horizon, consistent with the targeted
        holding period.
  \item Trade Intensity turns \textit{negative} at 15 minutes --- a warning sign
        discussed in Section~\ref{sec:discussion}.
  \item RSI and MA Crossover are the strongest individual signals at all horizons.
  \item The composite slightly outperforms the best individual signals at the
        1-minute horizon, reflecting the benefit of diversification across signal families.
\end{itemize}

\subsection{Quintile Decomposition}

Sorting bars into quintiles by signal strength and measuring the mean
5-minute forward return in each bin confirms a broadly monotonic relationship
for all price-based signals.  Bins 1 (most negative signal) and 5 (most positive)
show opposite-signed average returns in the range of $\pm$0.5--1 bps.  Trade
Intensity shows weak or reversed monotonicity, consistent with its low IC.

\subsection{Gross vs Net Performance: Continuous Composite}

\begin{center}
{\footnotesize
\begin{tabular}{lrrrrrr}
\toprule
\textbf{Signal} & \textbf{Gross PnL}
  & \textbf{Gross SR} & \textbf{Net PnL (2 bps)}
  & \textbf{Net SR (2 bps)} & \textbf{Turnover} & \textbf{Win\%} \\
\midrule
Z-score MR     & 0.031 & 34.9 & $-$0.111 & $-$118.9 & 708 & 51.9 \\
Bollinger MR   & 0.027 & 40.4 & $-$0.098 & $-$143.1 & 623 & 51.9 \\
VWAP MR        & 0.034 & 39.7 & $-$0.130 & $-$138.5 & 820 & 52.2 \\
OFI            & 0.045 & 56.1 & $-$0.152 & $-$178.0 & 982 & 50.4 \\
Trade Intensity& 0.001 &  1.4 & $-$0.424 & $-$505.4 & 2122 & 51.3 \\
RSI MR         & 0.036 & 49.1 & $-$0.117 & $-$150.3 & 767 & 52.3 \\
MA Crossover   & 0.026 & 39.8 & $-$0.049 & $-$74.9  & 373 & 52.9 \\
Vol. Momentum  & 0.026 & 32.5 & $-$0.096 & $-$113.5 & 609 & 51.4 \\
\midrule
\textbf{Composite}& \textbf{0.030} & \textbf{43.2}
  & $-$0.044 & $-$63.2 & 371 & 51.8 \\
\bottomrule
\end{tabular}
} % end footnotesize
\end{center}

Annualized Sharpe ratios assume 1-minute bars: $\sqrt{1440 \times 365}$ scaling.
All strategies go to zero or negative net PnL at the representative 2 bps
taker cost.  The composite is the most \textit{cost-efficient} signal: it
has the \textit{lowest turnover} among all strategies (371), a direct result
of the EMA smoothing that filters out high-frequency signal reversals.

\subsection{Transaction Cost Sensitivity: Composite Signal}

\begin{center}
\begin{tabular}{rrr}
\toprule
\textbf{Cost (bps)} & \textbf{Total PnL} & \textbf{Annualized Sharpe} \\
\midrule
0.00 & $+$0.030 & $+$43.2 \\
0.25 & $+$0.021 & $+$30.3 \\
0.50 & $+$0.012 & $+$17.5 \\
0.75 & $+$0.003 & $+$4.7  \\
1.00 & $-$0.006 & $-$8.0  \\
1.50 & $-$0.024 & $-$34.3 \\
2.00 & $-$0.044 & $-$63.2 \\
3.00 & $-$0.083 & $-$120.6 \\
5.00 & $-$0.162 & $-$235.3 \\
\bottomrule
\end{tabular}
\end{center}

The \textbf{breakeven transaction cost} is approximately \textbf{0.75 bps
round-trip} for the continuous composite signal.  This is a demanding
threshold: it requires limit-order (maker) execution with minimal adverse
selection, accessible only on exchange VIP fee tiers.

\subsection{Thresholded (Discrete) Position Strategy}

Using the hysteresis discretization from Section~\ref{sec:composite}:

\begin{center}
\begin{tabular}{rrrrr}
\toprule
\textbf{Cost (bps)} & \textbf{Total PnL} & \textbf{Annualized Sharpe}
  & \textbf{Profit Factor} & \textbf{Win\%} \\
\midrule
0.00 & $+$0.033 & $+$33.1 & 1.146 & 52.5 \\
0.25 & $+$0.023 & $+$23.3 & 1.099 & 48.8 \\
0.50 & $+$0.014 & $+$13.5 & 1.056 & 48.4 \\
0.75 & $+$0.004 & $+$3.8  & 1.015 & 48.1 \\
1.00 & $-$0.006 & $-$5.9  & 0.977 & 47.8 \\
2.00 & $-$0.045 & $-$42.8 & 0.847 & 47.0 \\
\bottomrule
\end{tabular}
\end{center}

The thresholded strategy produces \textbf{351 position changes} (175 round
trips) over 2 days with a mean holding period of 14.7 bars, which falls squarely
in the 5--15 minute target window.  Its lower gross Sharpe (33 vs 43) relative
to the continuous strategy reflects the information lost in discretization,
but its \textit{total turnover} (391 vs 371) is similar, since the hysteresis
suppresses trivial trades.

\subsection{Rebalance Frequency Sweep}

Holding the composite signal fixed and varying the rebalance interval
$\Delta$ (how often the position is updated):

\begin{center}
\begin{tabular}{lrrrrr}
\toprule
 & \multicolumn{5}{c}{\textbf{Annualized Sharpe by Cost (bps)}} \\
\cmidrule(lr){2-6}
\textbf{Rebal (min)} & \textbf{0} & \textbf{0.5} & \textbf{1.0} & \textbf{2.0} \\
\midrule
1  & 43.2 & 17.5 & $-$8.0  & $-$63.2 \\
3  & 35.1 & 18.4 & $+$1.8  & $-$43.8 \\
5  & 28.7 & 16.2 & $+$3.6  & $-$33.3 \\
7  & 24.1 & 13.5 & $+$2.9  & $-$26.5 \\
10 & 18.4 &  9.2 & $+$0.0  & $-$19.1 \\
15 &  9.6 &  2.7 & $-$4.1  & $-$16.8 \\
\bottomrule
\end{tabular}
\end{center}

Gross Sharpe declines monotonically with rebalance interval: the signal
edge is concentrated at the 1--3 minute horizon and has largely decayed
by 15 minutes.  For a cost of 0.5 bps, the peak risk-adjusted return is
achieved at a 3-minute rebalance, not the longer intervals one might naively
expect from a ``5--15 minute strategy''.

\subsection{Walk-Forward Out-of-Sample Results}

\begin{center}
\begin{tabular}{lrr}
\toprule
\textbf{Metric} & \textbf{Gross} & \textbf{0.5 bps} \\
\midrule
Total PnL                 & $+$0.025 & $+$0.011 \\
Annualized Sharpe         & $+$47.2  & $+$20.1  \\
Max drawdown              & $-$0.004 & $-$0.004 \\
Win rate                  & 51.4\%   & 41.8\%   \\
Profit factor             & 1.25     & 1.10     \\
Bars covered (OOS only)   & 2{,}160  & 2{,}160  \\
\bottomrule
\end{tabular}
\end{center}

The walk-forward Sharpe (47.2 gross) slightly \textit{exceeds} the in-sample
result (43.2), which may seem surprising.  The most likely explanation is
that the walk-forward drops the first 720 bars of warm-up period (where
signals are noisier) from the PnL computation, leaving a cleaner signal in
the test folds.

%% -----------------------------------------------------------------------
\section{Critical Discussion}
\label{sec:discussion}

\subsection{Signal Redundancy: Z-Score and Bollinger}

Signals 1 and 2 are essentially the same transformation.  Both compute
$(C_t - \bar{C}) / \sigma_C$; the Bollinger version merely clips the result
to $[-1, 1]$ and applies a fixed scaling by $k=2$.  Their empirical IC values
are identical to three decimal places (0.0463 vs 0.0462 at the 1-min horizon).
In a production signal suite, these would be merged into a single entry and
the freed weight reallocated to a genuinely orthogonal source of alpha ---
for example, funding-rate deviation, open-interest change, or a cross-instrument
(ETH) signal.

\subsection{Trade Intensity is Mis-Specified}

Signal 5 is designed as a momentum signal (high volume + up move = long)
inside a predominantly contrarian framework.  Its IC at the 15-minute horizon
is \textbf{negative} ($-0.018$): the strategy would be better served by
\textit{fading} high-volume moves, consistent with the other signals.
More fundamentally, trade intensity measured from 1-minute bars captures
aggregate activity rather than the order-by-order dynamics that make activity
measures genuinely informative.  A tick-level measure (e.g.\ consecutive same-side
trades, or the Lee-Ready-adjusted trade count) would be more appropriate.

\subsection{The Walk-Forward Does Not Test Parameter Stability}

The walk-forward scheme trains over a 12-hour window and tests over a 4-hour
window, but \textit{signal parameters are never re-estimated}.  The 20-bar
lookback for z-score, the 0.15 entry threshold for hysteresis, and the
signal weights are identical across all folds.  The walk-forward therefore
tests \textit{temporal stability of performance under fixed parameters}, not
the ability to generalize a fitting procedure.  A proper walk-forward would
optimize the parameters on each training window and evaluate them on the
subsequent test window, which would require a much larger dataset.

\subsection{The Sharpe Numbers are Overstated}

Annualized Sharpe ratios in the range of 30--50 should not be taken at face value.
Several effects combine to inflate them:

\begin{enumerate}[leftmargin=2em, itemsep=0pt]
  \item \textbf{One-bar delay is insufficient for HFT mean reversion.}
        At the 1-minute bar level, a mean-reversion strategy needs to submit
        limit orders \textit{within} the bar that experiences the dislocation,
        not at the close of the next bar.  The one-bar-later entry captures
        only the portion of the reversion that extends into subsequent bars.
  \item \textbf{Zero market impact.}  Any realistically sized position in BTCUSDT
        perps would move the price, particularly during low-volume periods.
        The backtest treats every fill as if it were executed at the mid with
        zero price impact.
  \item \textbf{Two days of data.}  The statistical uncertainty on a Sharpe
        ratio estimated from 2{,}880 bars is enormous.
        A 95\% confidence interval for a true Sharpe of 2 (annualized),
        derived from one quarter of data, can easily span $\pm 1.5$.
        Two days provides far less.
  \item \textbf{No adverse selection model for maker fills.}
        Limit orders at the bid fill more often when large sellers are
        incoming, which is the worst possible time to be long.
        Ignoring adverse selection systematically inflates maker strategy PnL.
\end{enumerate}

\subsection{The 0.75 bps Breakeven is an Optimistic Bound}

On Binance, VIP maker rebates can reach $-0.01$ bps (i.e.\ the exchange
\textit{pays} you for liquidity), which on paper makes the gross alpha
sufficient.  In practice:

\begin{itemize}[leftmargin=2em, itemsep=2pt]
  \item \textbf{Fill rate uncertainty.}  Maker orders are not guaranteed to fill.
        If the fill rate is 60\%, the effective cost per intended-trade is
        higher than the quoted fee tier.
  \item \textbf{Queue position.}  In a fast-moving order book, a new limit
        order joins at the back of the queue.  For a mean-reversion trade,
        where price may start recovering before the order reaches the front,
        fill quality degrades.
  \item \textbf{Co-location and infrastructure.}  Low-latency execution
        systems (co-location, direct market access, kernel bypass networking)
        have costs that should be amortized against strategy PnL.
\end{itemize}

\subsection{Two Days is Not a Dataset}

The most fundamental limitation is sample size.  Two trading days capture
a single, relatively quiet market regime (BTC price ranging \$84k--\$86k,
a 1.5\% window).  There is no ability to test:
\begin{itemize}[leftmargin=2em, itemsep=0pt]
  \item How the signal performs in high-volatility regimes (post-news, liquidation
        cascades, funding spikes).
  \item Whether the negative lag-1 autocorrelation is a stable feature or a
        regime-specific artefact.
  \item Seasonal effects (weekend vs weekday, options expiry days).
  \item Longer drawdown periods that would expose the fragility of the
        hysteresis thresholds.
\end{itemize}

\subsection{What Would Make This Credible}

For the strategy to be taken seriously in a production context, the following
would need to be demonstrated:

\begin{enumerate}[leftmargin=2em, itemsep=2pt]
  \item \textbf{$\geq 6$ months of tick data} spanning multiple volatility regimes.
  \item \textbf{Proper walk-forward with parameter optimization}: LASSO/ridge
        regression on signal weights, or a simple equal-weight ensemble baseline,
        estimated in-sample and frozen out-of-sample.
  \item \textbf{Queue simulation}: modeling fill probability as a function of
        order size, time-in-queue, and market depth.
  \item \textbf{IC t-statistics with multiple-testing correction} across the
        8 signals to assess which ones are genuinely informative.
  \item \textbf{Intraday volume normalization}: scaling signals by the expected
        volume in a given hour to avoid regime-mixing.
\end{enumerate}

%% -----------------------------------------------------------------------
\section{Conclusion}

This project demonstrates a complete quantitative research pipeline from raw
tick data to cost-aware backtesting for BTCUSDT perpetual futures.
The key empirical finding --- \textbf{negative lag-1 autocorrelation in 1-minute
returns} --- provides a sound microstructure basis for a mean-reversion signal
suite.  All price-based signals show positive information coefficient at the
5--10 minute horizon, and the composite signal achieves an IC of approximately
0.10, which is a reasonable starting point for a high-frequency mean-reversion
strategy.

The backtesting results indicate \textbf{gross alpha that is real but thin}.
The strategy is profitable in simulation only below $\sim$0.75 bps in round-trip
transaction costs --- a threshold achievable with maker limit-order execution
on Binance VIP tiers, but not with standard taker fills.

The critical analysis in Section~\ref{sec:discussion} highlights several areas
where the simulation is optimistic: signal redundancy between z-score and
Bollinger signals, a mis-specified momentum signal (Trade Intensity), inflated
Sharpe ratios due to unrealistic fill assumptions, and the fundamental limitation
of a two-day dataset.

The work should be viewed as a \textbf{research proof-of-concept} rather than
a deployable strategy: it identifies a plausible alpha source, validates it
with appropriate IC and quintile analysis, and quantifies the execution-cost
constraint clearly.  Extending the dataset and incorporating realistic execution
simulation would be the natural next steps before live deployment.

\end{document}
